\chapter{Context Graphs and Cycle Consistency}
\label{ch:context-graphs}

Chapter~\ref{ch:composition} asked whether projecting before composing
preserves the composition law for a fixed pair of lifts. This chapter asks a
different question entirely: given several \emph{independently observed}
transitions between measurement contexts, can they all have come from one
consistent global assignment of frames? The first question is about a
single, chosen composition; the second is about whether many separately
admissible observations are mutually compatible at all. Nothing about
Chapter~\ref{ch:composition}'s obstruction is assumed or needed here.

\section{Measurement Contexts as Vertices, Transitions as Edges}

Let $G = (V, E)$ be a graph whose vertices are measurement contexts (orthonormal
bases) and whose edges carry observed transition-probability matrices $B_e$. We
distinguish four increasingly demanding notions of admissibility over a field
$\mathbb F \in \{\R, \C, \Hset\}$:
\begin{center}
\begin{tabular}{@{}ll@{}}
\toprule
vertex data & a context, prior to any transition data \\
edge admissibility & $B_e$ individually lies in the $\mathbb F$-admissible set \\
path composition & sequential lifts compose (Chapter~\ref{ch:composition}) \\
cycle consistency & lifts around a closed loop compose to the identity \\
global realizability & a single assignment of frames realizes every edge at once \\
\bottomrule
\end{tabular}
\end{center}

\begin{definition}[Edgewise vs.\ global $\mathbb F$-admissibility]
Write $\mathcal A_{\mathbb F}(e)$ for the set of edges individually admitting an
$\mathbb F$-lift, and $\mathcal A_{\mathbb F}(G)$ for the set of context graphs
admitting a \emph{single} assignment of frames to vertices realizing every edge
simultaneously.
\end{definition}

\section{Exact Closure and Its Failure Under Projection}

For exact frame-to-frame transformations $L_{AB}$ around a triangle of contexts
$Z, X, Y$,
\begin{equation}
L_{ZX} L_{XY} L_{YZ} = I.
\label{eq:exact-closure}
\end{equation}
Applying the probability projection $\pi(L) = |L|^{\circ 2}$ destroys the
information needed to express the corresponding closure condition: there is in
general no relation
\begin{equation}
B_{ZX} B_{XY} B_{YZ} = I
\end{equation}
satisfied by the observed transition matrices alone.

\begin{ledgerentry}
Equation~\eqref{eq:exact-closure} projecting to no corresponding closure relation
on the observed $B$'s is the central fact of this chapter: exact composition has
an identity that probability-level projection does not preserve.
\end{ledgerentry}

\section{Gauge Bookkeeping for a Single Real Frame}
\label{sec:gauge-single-frame}

Before turning to cycles and global consistency, it is worth asking a more
basic question about a single edge: recall $\theta \in [0,\pi/4]$ from
Chapter~\ref{ch:unistochastic}, the canonical unsigned transition angle
between two real two-dimensional contexts. Two very different things
compress a raw frame angle $\alpha$ down to this domain, and confusing them
would understate what $\pi$ (the probability projection) actually costs: one
is harmless relabeling, the other is a genuine, single-edge loss of
information that happens \emph{before} any cycle or composition enters the
picture. Separating them precisely is the point of this section.

\begin{proposition}[Gauge versus projection loss in the fundamental domain]
\label{prop:gauge-vs-loss}
\index{Gauge versus projection loss in the fundamental domain (proposition)}
The canonical unsigned transition angle $\theta \in [0,\pi/4]$ arises from composing
two identifications of \emph{different character}, not one redundancy in
disguise:
\begin{enumerate}
\item[(i)] \textbf{Gauge (label swap).} $\alpha \sim \alpha + \pi/2$. This
identifies the \emph{same} pair of lines with the labels "$+$" and "$-$"
exchanged: explicitly $e_+(\alpha+\pi/2) = e_-(\alpha)$ up to sign. Nothing
physical changes; only $B_{11} \mapsto 1 - B_{11}$.
\item[(ii)] \textbf{Projection loss, not gauge.} $\alpha \sim \pi - \alpha$.
Here $e_+(\pi-\alpha) \neq \pm e_+(\alpha)$ for generic $\alpha$: these are
\emph{genuinely different} lines. Yet $B_{11}(\alpha) = B_{11}(\pi-\alpha)$
identically, since $\cos^2\alpha = \cos^2(\pi-\alpha)$. The map $\pi$ destroys
the distinction between these two physically distinct contexts already at the
level of a single edge, prior to any composition or cycle.
\end{enumerate}
\end{proposition}

\begin{proof}[Verification]
Take $\alpha = \pi/6$. Direct computation gives $B_{11}(\pi/6) = 3/4$,
$B_{11}(\pi/6+\pi/2) = 1/4$ (confirming (i), $t \mapsto 1-t$), and
$B_{11}(\pi-\pi/6) = 3/4$ while $e_+(\pi-\pi/6) \neq \pm e_+(\pi/6)$ (confirming
(ii): identical observable, distinct context). Composing (i) and (ii) generates
the four generic images $\{\alpha,\ \pi/2+\alpha,\ \pi-\alpha,\ \pi/2-\alpha\}$
of a representative $\alpha$, an order-4 group consistent with the $[0,\pi/4]$
fundamental domain.
\end{proof}

\begin{remark}
This confirms the framing above precisely: part (ii) is a \emph{single}-edge
loss of information under $\pi$, independent of anything about cycles or
global consistency. The cycle-consistency obstruction of
\S\ref{ch:real-mub-obstruction} is therefore an \emph{additional},
qualitatively different failure (inability to glue) stacked on top of this more
basic single-edge failure (inability to invert $\pi$). The two should not be
conflated: the $\pm$ signs in the triangle-gluing relation
(Theorem~\ref{thm:triangle-gluing}) track both ambiguities at once, which is
why four sign combinations appear rather than two.
\end{remark}

\section{Global Gauge for a Context Graph}

A gauge decision at one vertex acts on every edge incident to it: for allowed
representative changes $g_A$ at each vertex,
\begin{equation}
L_{AB} \mapsto g_A^{-1} L_{AB} g_B.
\end{equation}
Consequently a triangle of contexts is not three independent two-basis problems;
signs cannot be chosen edge by edge. Fixing $\alpha_Z = 0$, $\alpha_X = \alpha$,
$\alpha_Y = \beta$, the exact relative angles obey
\begin{equation}
\Delta_{ZX} + \Delta_{XY} + \Delta_{YZ} = 0 \pmod{2\pi},
\end{equation}
which projects, after probability observation $\Delta \mapsto \cos^2\Delta$, to
the weaker triangle-gluing relation of the next section.

\section{Spanning Trees, Cycle Rank, and General Realizability}
\label{sec:spanning-trees}

The triangle of Chapter~\ref{ch:real-mub-obstruction} is the smallest instance
of a general phenomenon, worth isolating for its own sake before specializing
back to $n=3$ contexts.

\begin{definition}[Context graph]
A \emph{context graph} is a connected graph $G=(V,E)$ whose vertices are
measurement contexts and whose edges $e=(A,B)$ each carry an observed
bistochastic matrix $B_e$. Its \emph{cycle rank} (first Betti number) is
\begin{equation}
\beta_1(G) = |E| - |V| + 1.
\end{equation}
\end{definition}

A spanning tree $T \subseteq E$ has $|T| = |V|-1$ edges and $\beta_1(T) = 0$;
each additional edge beyond a spanning tree closes exactly one independent
cycle, so $\beta_1(G)$ counts the non-tree edges of any fixed spanning tree.

\begin{proposition}[Trees carry no independent constraint]
\label{prop:tree-no-constraint}
\index{Trees carry no independent constraint (proposition)}
Let $T=(V,E_T)$ be a tree, each edge individually $\mathbb F$-admissible (some
lift exists reproducing $B_e$). Then a global frame assignment reproducing
every edge of $T$ simultaneously always exists.
\end{proposition}

\begin{proof}
Fix a root $r \in V$ and $F_r := I$. $T$ has a unique path from $r$ to any
other vertex $v$; choose, for each tree edge, any one admissible lift, and
propagate $F_v := F_u L_{uv}$ along that unique path. Since the path to each
vertex is unique (this is exactly what makes $T$ a tree), no two propagation
routes can disagree, and by construction $F_A^{-1}F_B = L_{AB}$ reproduces
$B_e$ on every tree edge.
\end{proof}

\begin{theorem}[Graph realizability criterion]
\label{thm:graph-realizability}
\index{Graph realizability criterion (theorem)}
Let $G=(V,E)$ be a context graph and $T$ a spanning tree. $G$ is globally
$\mathbb F$-realizable (some assignment of vertex frames reproduces every edge
of $E$) if and only if there is a choice of admissible lifts on the tree edges
of $T$ such that the resulting frame assignment (Proposition
\ref{prop:tree-no-constraint}) \emph{also} reproduces $B_e$ on every non-tree
edge $e \in E \setminus T$.

Equivalently: global realizability holds iff there exists a \emph{single}
compatible choice of lifts on a spanning tree for which \emph{every} non-tree
edge simultaneously satisfies its induced fundamental-cycle closure condition.
The existential quantifier ranges over one shared tree-lift choice, not over
each non-tree edge separately — this is the content of the theorem, not a
footnote to it.
\end{theorem}

\begin{proof}
($\Leftarrow$) The constructed frame assignment reproduces every tree edge by
Proposition~\ref{prop:tree-no-constraint} and every non-tree edge by
hypothesis, hence every edge of $E$.

($\Rightarrow$) If some global frame assignment $\{F_A\}$ reproduces every
edge of $E$, its restriction to $T$ is in particular a choice of admissible
tree-edge lifts ($L_{AB} := F_A^{-1}F_B$ for $(A,B) \in T$ reproduces $B_e$ by
hypothesis), and propagating from that choice reconstructs exactly $\{F_A\}$
(by uniqueness of tree paths), which by hypothesis already matches every
non-tree edge.
\end{proof}

The theorem is existential rather than constructive: it characterizes
\emph{when} a compatible assignment exists without prescribing a unique one
or a search procedure for finding it. Multiple tree-lift choices can succeed
simultaneously (Chapter~\ref{ch:real-mub-obstruction}'s explicit real
solutions to the triangle-gluing relation are one instance), and nothing
here selects among them.

\begin{remark}
Theorem~\ref{thm:graph-realizability} reduces global realizability to exactly
$\beta_1(G)$ simultaneous consistency checks — one per non-tree edge, all
against the \emph{same} tree-lift choice, not checked independently of one
another. This is a genuine constraint when $\beta_1(G) \ge 2$: passing each
non-tree check in isolation, for possibly different tree-lift choices, does
not suffice: it must be a single choice of tree lifts that satisfies every
non-tree edge at once.
\end{remark}

\begin{corollary}
\label{cor:triangle-minimal}
$\beta_1(K_3) = 3-3+1 = 1$: the triangle of
Chapter~\ref{ch:real-mub-obstruction} is the minimal context graph possessing
a nontrivial cycle constraint, and Theorem~\ref{thm:triangle-gluing} is the
$\beta_1=1$ case of Theorem~\ref{thm:graph-realizability} (spanning tree
$\{ZX,ZY\}$, single non-tree edge $XY$).
\end{corollary}

\begin{remark}[Two independent axes, not one ladder]
\label{rem:two-axes}
Theorem~\ref{thm:graph-realizability} should not be read as one more rung on a
single vertex$\to$edge$\to$path$\to$cycle$\to$global hierarchy. Part I
actually establishes two \emph{logically independent} questions, and
collapsing them into one ordered scale obscures that independence.

\emph{Compositional sufficiency} (Chapter~\ref{ch:composition}) asks whether
projected data retains enough information to predict sequential composition:
$\pi(L_2L_1) \overset{?}{=} \pi(L_2)\pi(L_1)$. This question is already
meaningful on a bare open path — no cycle, and in particular no closed loop
in a context graph, is required to ask or answer it. It measures information
lost by projection, full stop.

\emph{Global realizability} (this chapter and Chapter~\ref{ch:real-mub-obstruction})
asks whether independently specified, individually admissible edge data
$\{B_e\}$ on a context graph admits one simultaneous vertex-frame assignment
reproducing all of it at once. This is where trees, non-tree edges, and
$\beta_1(G)$ belong, and it is silent on whether any particular composition of
chosen lifts is faithful.

The MUB triangle fails the second test; the interference calculation of
Chapter~\ref{ch:composition} demonstrates failure of the first. Nothing forces
these into a common scale. An assumption's contribution to the reconstruction
is better recorded as a pair
\begin{equation}
Q(A) = \big(Q_{\mathrm{comp}}(A),\ Q_{\mathrm{glob}}(A)\big)
\end{equation}
— what it contributes to compositional recovery, and separately what it
contributes to global realizability — rather than a single numerical "depth."
Part II introduces a third, independent axis in the same spirit: \emph{system
composition} (does an individually admissible state space for each of several
systems compose into an admissible joint state space?), which is neither
sequential composition along a path nor contextual gluing around a cycle.
\end{remark}
